% !PRINTED 254
\weberpage{254}

Now if $x_1$ is a real root of (1), then $y_1^2$ is a positive square in the sum (4). But if $x_1$ and $x_2$ are conjugate imaginary, then $y_1$ and $y_2$ likewise split into two conjugate imaginary constituents
\[
y_1 = u + vi, \quad y_2 = u - vi,
\]
where $u$ and $v$ are real linear functions of the $t$. Accordingly we have
\[
y_1^2 + y_2^2 = 2u^2 - 2v^2.
\]

% The source visibly letterspaces only this occurrence of "verschiedenen Wurzeln".
Thus by (4) the function $\Phi(t)$ defined by (5) is decomposed into a sum of squares, whose number is equal to the number of \gesperrt{distinct roots} of $f(x) = 0$, and among which there are as many negative ones as there are pairs of imaginary roots among these mutually distinct roots.

But by the law of inertia of quadratic forms (\S.\,58) these numbers remain the same under every other real linear transformation of the quadratic form into a sum of squares. Now $S(y)$ is a real linear function of the $t$, and $B$ no longer contains the variable $t_{n-1}$; from this, formula (5) yields the theorem:

\gesperrt{I. If the Bezoutiant $B$ is decomposed by a real linear transformation into a sum of $\pi$ positive and $\nu$ negative squares, which cannot be reduced to a smaller number, then $\pi + \nu + 1$ is the number of distinct roots of $f(x) = 0$; and $\nu$ is the number of pairs of conjugate imaginary roots contained among them, $\pi - \nu + 1$ that of the real ones.}

If the determinant of $B$, and hence also the discriminant of $f(x)$ (\S.\,73), is different from zero, then $\pi + \nu = n - 1$ and $\nu$ is less than or at most equal to $\frac{1}{2}n$.

The Bezoutiant is thus a quadratic form of a \gesperrt{special nature}, which expresses itself precisely in the fact that among the squares into which it can be decomposed, at most $\frac{1}{2}n$ negative ones can occur.
